\documentclass[12pt]{article}
\usepackage{graphicx} % for creating figures
\graphicspath{ {figures/} }

\usepackage{amsmath} % for math environments
\usepackage{amsfonts} % for math fonts like \mathbb
\usepackage{cancel} % for striking out math
\usepackage{bm} % for bold math



% stop typing \mathbb a thousand times 
\newcommand{\R}{\mathbb{R}}
\newcommand{\C}{\mathbb{C}}
\newcommand{\F}{\mathbb{F}}
\newcommand{\E}{\mathbb{E}}
\newcommand{\M}{\mathbb{M}}
\newcommand{\Z}{\mathbb{Z}}
\newcommand{\sphere}{\mathbb{S}}

\newcommand{\defeq}{\vcentcolon=} % for a prettier := sign
\newcommand{\eqdef}{=\vcentcolon} % for a prettier =: sign


% commands for bra-ket notation
\newcommand{\bra}[1]{\left\langle#1\right|}
\newcommand{\ket}[1]{\left|#1\right\rangle}
\newcommand{\braket}[2]{\left\langle #1 \middle| #2 \right\rangle}
\newcommand{\matrixel}[3]{\left\langle #1 \middle| #2 \middle| #3 \right\rangle}
\newcommand{\expectation}[1]{\left\langle #1 \right\rangle}


% vector stuff
\newcommand{\basis}{\bm{\hat{e}}}
\newcommand{\unit}[1]{\bm{\hat{#1}}}
\newcommand{\bvec}[1]{\bm{\vec{#1}}}
\newcommand{\del}{\bm{\vec{\nabla}}}

% curly velocity v 
\newcommand{\velocity}[1]{\vec{\mathcal{v}}_{{}_{#1}}}



\title{Curves in $\R^3$}

\begin{document}
\maketitle
In order to discuss the theory of curves in $\R^3$ we must establish some
preliminary definitions.

\textbf{Definition}(\textit{smooth}): A (real) function is said to be \textit{smooth} if for every
point it has continuous derivatives of any order.

With this definition, a parametrized curve is just a smooth map from some
interval of the reals to the space in question. That is,
\textbf{Definition}(\textit{parametrized curve}): 
A \textit{parametrized curve} is a smooth map
\begin{equation}
  \vec{r}:(a,b)\subseteq \R \to \R^3
\end{equation}
where
\begin{equation}
  \vec{r}(\lambda): t\mapsto (x(\lambda),\; y(\lambda),\; z(\lambda)).
\end{equation}
In other words, the coordinate functions $x(\lambda),\;y(\lambda),\;z(\lambda)$ are smooth. The
variable $\lambda$ is called the \textit{parameter} of the curve. Note that we
have chosen to use $\lambda$ instead of $t$ because the parameter need not be
``time''. This becomes very important in relativity where time is treated as a
fourth dimension instead of a universal parameter. In this instance, 

Consequently, we will not call the derivative of the curve with respect to our
parameter ``the velocity'' and instead will use \textit{tangent vector}. That
is,

\textbf{Definition}(\textit{tangent vector}): The derivative of a paramterized
curve $\vec{r}(\lambda)$ with respect to the parameter $\lambda$ at $\lambda = \lambda_0$
\begin{equation}
  \left.\frac{d\vec{r}}{d\lambda}\right|_{\lambda_0} = (x'(\lambda_0),\; y'(\lambda_0),\; z'(\lambda_0))
\end{equation}
is called the \textit{tangent vector} to the curve $\vec{r}(\lambda)$ at the
point $\lambda=\lambda_0$.

What properties of curves should we care about? A perhaps more useful
reformulation is \textit{What properties can we use to compare two
  curves?} Everyday experience leads to two key concepts:
\begin{itemize}
\item The length of the curve (arclength)
\item The ``bendiness'' of the curve (curvature)
\end{itemize}

Let's start with the length of the curve. At every point along the curve
$\vec{r}(\lambda)$ we can define a tangent vector
\begin{equation}
  \velocity{\lambda_0} = \left.\frac{d\vec{r}(\lambda)}{d\lambda}\right|_{\lambda_0}
\end{equation}
The span of our tangent vector $\velocity{\lambda_0}$ defines the tangent line
to $\vec{r}$ at $\lambda_0$. Therefore, we can consider partitioning our curve
into segments such that $\lambda\in \left\{\lambda_i\right\}_{i=0}^{n}$ and
approximating the curve by piecewise connected tangent lines. The length of the
approximated curve is the sum of the lengths of each tangent line
$\sum|\velocity{\lambda_i}|\Delta \lambda$. Thus, in the appropriate limit, we have
\begin{equation}
  s(\lambda) = \int_{\lambda_0}^{\lambda}|\velocity{\lambda}|d\lambda
\end{equation}
which we take as \textbf{the definition of the arclength of $\vec{r}(\lambda)$}.
Thus, we have reduced calculating the lengths of curves to calculating the
lengths of tangent vectors. 

To compute $|\velocity{\lambda}|$, we observe that $\velocity{\lambda}\in\R^3$.
Thinking of $\R^3$ as a vector space, we define the length of a tangent
vector by the Euclidean inner product
\begin{equation}
  |\velocity{\lambda}| = \sqrt{\velocity{\lambda}\cdot\velocity{\lambda}} = \sqrt{\sum_i v_i^2}.
\end{equation}
\textbf{Note:} We chose to use the Euclidean inner product because it aligns
with our everyday sense of length. This is just a choice! If evidence dictates
that a different inner product should be used, we can pick something else.

\textbf{Note 2:} If the parameter of our curve is physical time $t$, then this definition says
that the length of a curve is just the integral of the speed with respect to
time.

What if we had chosen to parametrize our curve by the arclength. In other words,
$\lambda = s$? Then
\begin{align}
  \frac{ds}{d\lambda} &= |\velocity{\lambda}| = 1  \\
  \Rightarrow s &= \int_{\lambda_0}^{\lambda}d\lambda = \lambda-\lambda_0
\end{align}

\end{document}






























